% chapter1.tex
\chapter{کلیات تحقیق}

\section{مقدمه}
(در این بخش پژوهشگر به ارائه اطلاعات دقیق و روشن درباره موضوع مورد پژوهش، ارائه دلایل منطقی و نظری منجر به انجام پژوهش و بیان ارتباط بین این پژوهش با پژوهش‌های قبلی می‌پردازد.)

\section{بیان مسئله}
(در این بخش پژوهشگر به بیان ابهام‌ها، چالش‌ها، شکاف‌های دانشی، تعارض بین داده‌های پیشین، موارد مجهول و نیازهای موجود در رابطه با موضوع پژوهش می‌پردازد.)

\section{اهمیت و ضرورت پژوهش}
(در این بخش پژوهشگر باید درباره اهمیت، مزایا و اولویت‌های انجام پژوهش توضیح داده و ضرورت و تبعات ناشی از عدم انجام این پژوهش را بیان کند.)

\section{هدف پژوهش}
(اهداف پژوهش بیان نظام‌مند مواردی است که پژوهشگر با توجه به مسئله پژوهش به دنبال دستیابی به آن است.)

\section{پرسش‌ها یا فرضیه‌های پژوهش}
(پرسش‌های پژوهش باید متناسب و کاملاً همخوان با اهداف تنظیم شوند. متناسب با موضوع پژوهش ممکن است به جای پرسش از فرضیه استفاده شود.)

\section{تعاریف مفهومی و عملیاتی اصطلاحات}
(در تعریف مفهومی، یک اصطلاح تخصصی با استناد به منابع معتبر تعریف می‌شود. در تعریف عملیاتی، ویژگی‌ها و شاخص‌های مورد استفاده در این پژوهش بیان می‌شود.)